%%%%%%%%%%%%%%%%%%%%%%%%%%%%%%%%%
%
%       EXERCÍCIO
%
%%%%%%%%%%%%%%%%%%%%%%%%%%%%%%%%%

\definevalorquestao

\begin{exercicioBanco}[\valorquestao]
Um retângulo inicial, de perímetro 200 centímetros, sofre uma modificação tal que a medida de sua largura aumenta 20\%, e a medida do seu comprimento diminui 20\%. Determine a função que define a área \(A\) do novo retângulo, em centímetros quadrados, em relação à medida da largura do retângulo inicial \(x\), em centímetros.

% Define as alternativas
\newcommand{\alternativas}{%
    \begin{center}
        \begin{tabularx}{\textwidth}{XXX}
            (a) \(A(x) = 120x - 0,8x^{2}\). &
            (b) \(A(x) = 120x + 0,8x^{2}\). &
            (c) \(A(x) = 98x - 0,98x^{2}\). \\[5pt]
            (d) \(A(x) = 80x - 1,2x^{2}\). &
            (e) \(A(x) = 96x - 0,96x^{2}\).
        \end{tabularx}
    \end{center}
}

% Define a resposta correta
\newcommand{\resposta}{E}

% Lógica condicional para exibição
\ifdefstring{\atividade}{lista}{%
    \alternativas
    \vspace{0.5em}
    
    \noindent\textbf{Resposta:} letra \textbf{\resposta}.
}{%
    \ifdefstring{\modo}{objetivo}{%
        \alternativas
    }{%
        % modo = discursiva → não mostra alternativas
    }
}
\end{exercicioBanco}

